\documentclass[10pt,a4paper]{article}

\usepackage[margin=20mm]{geometry}
\usepackage{amsmath,amssymb,amsthm}

\theoremstyle{definition}
\newtheorem*{definition*}{Definition}
\newtheorem*{problem*}{Problem}

\newcommand{\Lclass}{\mathsf{L}}
\newcommand{\BPL}{\mathsf{BPL}}
\newcommand{\BPSPACE}{\mathsf{BPSPACE}}
\newcommand{\DSPACE}{\mathsf{DSPACE}}

\setlength{\parindent}{0pt}
\setlength{\parskip}{5pt}
\setlength{\emergencystretch}{2em}

\begin{document}

\begin{center}
  {\Large\bfseries BPL versus Deterministic Logspace}
\end{center}

\section{Definitions}

\begin{definition*}[Deterministic logspace]
$\Lclass$ is the class of languages decided by deterministic Turing machines
using $O(\log n)$ cells of read-write work tape on inputs of length $n$. The
input tape is read only, and the machine halts on every input.
\end{definition*}

\begin{definition*}[Bounded-error probabilistic logspace]
A probabilistic logspace machine has a read-only input tape, an $O(\log n)$-cell
work tape, and a one-way read-once random tape containing independent unbiased
bits. A language $A$ is in $\BPL$ if there are such a machine $M$ and a
polynomial $p$ such that, for every $x\in\{0,1\}^n$ and every random string
$r\in\{0,1\}^{p(n)}$, the computation $M(x;r)$ halts within $p(n)$ steps and
outputs one bit, with
\begin{align}
  x\in A
  &\Longrightarrow
  \Pr_{r\leftarrow U_{p(n)}}[M(x;r)=1]\geq\frac{2}{3},\\
  x\notin A
  &\Longrightarrow
  \Pr_{r\leftarrow U_{p(n)}}[M(x;r)=1]\leq\frac{1}{3}.
\end{align}
The constants $2/3$ and $1/3$ can be replaced by any two fixed constants
separated from $1/2$: independent repetition and majority preserve logarithmic
space and polynomial time. Plainly,
\begin{equation}
  \Lclass\subseteq\BPL.
\end{equation}
More generally, $\BPSPACE(S)$ denotes the analogous bounded-error class with
$O(S(n))$ work space and with every random computation path required to halt;
the polynomial-time restriction is not imposed on this larger class.
\end{definition*}

\begin{definition*}[Read-once branching programs and generators]
A length-$T$, width-$W$ read-once branching program is a layered directed
graph with layers $V_0,\ldots,V_T$, at most $W$ vertices in each layer, a
start vertex in $V_0$, and two edges labelled $0$ and $1$ leaving every vertex
outside $V_T$. A string $r\in\{0,1\}^T$ selects one edge at each layer and
therefore determines an accept-reject output $B(r)\in\{0,1\}$.

A generator $G:\{0,1\}^d\to\{0,1\}^T$ $\varepsilon$-fools this class of
programs if every such $B$ satisfies
\begin{equation}
  \left|
    \Pr[B(U_T)=1]
    -
    \Pr[B(G(U_d))=1]
  \right|
  \leq\varepsilon.
\end{equation}
Fixing an input $x$ to a polynomial-time probabilistic logspace machine gives
a read-once branching program whose vertices are the machine configurations.
Thus, for input length $N$, both $T$ and $W$ are $N^{O(1)}$.
\end{definition*}

\newpage
\section{Best Known Result}

There is an explicit generator for length-$T$, width-$W$ read-once branching
programs with seed length
\begin{equation}
  O\left(
    \log T\cdot\log\frac{TW}{\varepsilon}
  \right).
\end{equation}
For $T,W=N^{O(1)}$ and constant error, this is $O(\log^2N)$. Enumerating all
seeds gives a deterministic $O(\log^2N)$-space simulation of $\BPL$. This
remains the best ordinary pseudorandom generator stated in the source for
general polynomial-length, polynomial-width branching programs in the
constant-error regime.

The landmark general simulation is
\begin{equation}
  \BPSPACE(S)
  \subseteq
  \DSPACE\bigl(O(S^{3/2})\bigr).
\end{equation}
Using weighted pseudorandom generators, the best general unconditional space
bound stated in the source is
\begin{equation}
  \BPSPACE(S)
  \subseteq
  \DSPACE\left(
    O\left(\frac{S^{3/2}}{\sqrt{\log S}}\right)
  \right),
\end{equation}
and consequently
\begin{equation}\label{eq:bpl-bound}
  \BPL
  \subseteq
  \DSPACE\left(
    O\left(\frac{\log^{3/2}N}{\sqrt{\log\log N}}\right)
  \right).
\end{equation}

A weighted generator estimates the uniform acceptance probability using a
short signed or weighted combination of pseudorandom outputs; unlike an
ordinary pseudorandom generator, it need not induce a genuine probability
distribution. Error reduction can achieve optimal dependence on the branching
program's alphabet size while matching the best other weighted-generator
parameters. This identifies deweightization as a route toward an ordinary
$o(\log^2N)$-seed generator, but it does not itself improve
the deterministic-space bound in~\eqref{eq:bpl-bound}.

\section{Problem Formalization}

\begin{problem*}[$\BPL$ versus $\Lclass$]
Prove
\begin{equation}
  \BPL=\Lclass.
\end{equation}
Major intermediate progress would include improving the general simulation
to
\begin{equation}
  \BPL
  \subseteq
  \DSPACE\left(
    o\left(\frac{\log^{3/2}N}{\sqrt{\log\log N}}\right)
  \right).
\end{equation}
\end{problem*}

\end{document}
